%%%%%%%%%%%%%%%%%%%%%%%%%%%%%%%%%%%%%%%%%%%%%%%%%%%%%%%%%%%%%%%%%%%%%%%%%%%%%%%%
%% 
\cleardoubleoddpage%  Make sure to start each chapter on a new odd page
\chapter{Discussion}
\label{sec:discussion}

This chapter discusses the experimental results in the context of biosignal processing research. It examines the performance improvements observed with \ac{xLSTM} variants and discusses what they mean for clinical applications. The discussion covers both the achievements and limitations of integrating \ac{xLSTM} architectures within the \ac{BIOT} framework.

\section{Performance Interpretation and Window Effects}
\label{sec:performance_analysis}

The experimental results show that performance depends on the window length used. There is no single architecture that performs best in all cases. \ac{sLSTM} achieves best performance at 5\,s ($52.15\% \pm 1.37\%$ balanced accuracy) and 9\,s windows ($52.46\% \pm 3.81\%$), but performs worse at 7\,s ($48.45\% \pm 5.50\%$). In contrast, \ac{mLSTM} performs best at 7\,s windows ($50.66\% \pm 2.87\%$) where \ac{sLSTM} performs worst. The Linear Transformer baseline performs well at 5\,s ($52.06\% \pm 4.06\%$) but gets worse with longer windows.

\subsection{Window-Dependent Benefits}
\label{sec:temporal_modeling}

The results show that \ac{xLSTM} does not consistently outperform the baseline. Instead, each memory mechanism works better with specific window lengths. At 5\,s windows, the Linear Transformer achieves 52.06\% balanced accuracy, which is almost the same as \ac{sLSTM}'s 52.15\%. The Linear Transformer also trains 2.3× faster. Only at 9\,s windows does \ac{sLSTM} show clear advantages ($52.46\%$ vs. $48.45\%$ for Linear Transformer), though this comes with longer training time (60\% slower) and high variance (±3.81\%).

Several factors can explain these limited and window-dependent benefits:

\textbf{Sequence-Length Mismatch:} The \ac{BIOT} tokenization creates relatively short sequences (9-17 tokens per channel for 5-9\,s windows, 144-272 total tokens). These sequences may not need the advanced long-term memory capabilities that \ac{xLSTM} provides. The enhanced memory structures designed for long-sequence language modeling may be too complex for biosignal sequences where patterns occur within shorter time spans.

\textbf{Local Pattern Sufficiency:} The good baseline performance suggests that local patterns captured by linear attention may be enough for seizure classification. Neurological events in the \ac{TUEV} dataset may be characterized by local features (frequency patterns, amplitude changes, channel correlations) that do not need complex long-range dependency modeling.

\textbf{Training Dynamics Dominance:} The high variability across runs (coefficient of variation 6-10\%) shows that training dynamics, random initialization, and optimization challenges have a bigger impact on performance than the architecture itself. This suggests that improving training stability may give better results than using more complex architectures for this task.

\textbf{Hyperparameter Suboptimality:} All architectures use hyperparameters optimized for the Linear Transformer baseline (dropout 0.2, 8 blocks, embedding dimension 256). The \ac{xLSTM} variants may not work well with these settings. The window-dependent patterns might show hyperparameter mismatch rather than true architectural limitations.

\subsection{Architecture-Specific Window Behaviors}
\label{sec:window_behaviors}

Each architecture responds differently to window length:

\textbf{Linear Transformer:} Shows consistent degradation with increasing window length (52.06\% → 50.08\% → 48.45\%). This suggests fundamental limitations in linear attention for longer sequences, where attention is spread across more tokens.

\textbf{\ac{sLSTM}:} Shows a U-shaped pattern with a strong drop at 7\,s windows. Class~2 recall drops from 52.5\% (5\,s) to 42.9\% (7\,s) with statistical significance (p=0.043). This window-specific failure may be caused by gradient flow issues at intermediate sequence lengths, memory capacity mismatch, or hyperparameter sensitivity. The recovery at 9\,s suggests scalar memory can adapt to longer sequences once they exceed a certain length.

\textbf{\ac{mLSTM}:} Shows an inverted pattern, peaking at 7\,s where \ac{sLSTM} fails but collapsing at 9\,s (46.40\%). This suggests different optimal operating ranges. The poor performance at 9\,s may be caused by higher memory requirements, difficulty maintaining key-value associations, or parallelization limits with very long sequences.

\textbf{Combined Architecture:} The \ac{mLSTM}+\ac{sLSTM} architecture shows very stable performance across windows ($50.03\%$ to $50.46\%$) but never achieves the best performance. The 7:1 ratio may reduce the benefits of each mechanism instead of combining them effectively.

\subsection{Domain Transfer Challenges}
\label{sec:domain_transfer}

The modest improvements contrast with \ac{xLSTM}'s large gains in natural language processing~\cite{beck2024xlstm}. Language has rich long-range dependencies requiring information from hundreds or thousands of tokens. In contrast, the tokenized \ac{EEG} sequences (9-17 tokens per channel) mainly encode local patterns. The \ac{BIOT} tokenization~\cite{yang2023biot} converts 0.5-second windows of frequency content into discrete tokens, creating sequences where relevant patterns may be captured within short time spans.

Additionally, biosignal tokens represent frequency-domain features from continuous signals without the discrete semantic structure of language~\cite{yang2023biot}. Neurological events in the \ac{TUEV} dataset~\cite{harati2015improved} may have more consistent temporal scales, reducing the value of \ac{xLSTM}'s adaptive memory mechanisms. The evidence suggests mainly computational rather than neurophysiological constraints drive the window-dependent patterns—each architecture shows different optima, suggesting patterns reflect architectural limitations rather than underlying neurophysiology.

\section{Statistical Significance and Clinical Limitations}
\label{sec:critical_assessment}

\subsection{Statistical Power and Sample Size}
\label{sec:statistical_significance}

The lack of statistical significance in most performance comparisons (p > 0.05) is an important limitation. Only two comparisons achieved significance: AUCPR Macro at 5\,s (p=0.005, t-test) and overall accuracy at 9\,s (p=0.029, t-test). Balanced accuracy—the primary metric—showed no significant differences. This suggests that training dynamics and optimization problems may be more important than architectural improvements for biosignal processing.

With only 5 runs per configuration, the study lacks enough statistical power to detect small but potentially meaningful differences. The large standard deviations (3-5\% for balanced accuracy) show that training dynamics, initialization, and dataset characteristics have a bigger impact on performance than architectural choices. The 4.4\% relative improvement (2.18 percentage points absolute) achieved by \ac{sLSTM}, while consistent, may not be clinically meaningful given overlapping confidence intervals.

\subsection{Class Imbalance and Deployment Barriers}
\label{sec:class_imbalance}

The confusion matrix analysis reveals severe limitations that question clinical viability:

\textbf{Severe Class Imbalance Impact:} Class~0, representing only 1.93\% of training data, never exceeds 15\% recall even in the best case (\ac{sLSTM} at 9\,s: 14.2\% ± 9.4\%). This means algorithms would miss approximately 85\% of minority-class events. In clinical applications, missing any seizure type could have severe patient safety implications. This makes them unsuitable for clinical decision support without additional strategies such as cost-sensitive learning~\cite{cui2019class}, focal loss~\cite{lin2017focal}, or targeted data augmentation.

\textbf{Majority Class Bias:} The 87.87\% accuracy on Class 5 (66.77\% of data) versus 0\% accuracy on Class 0 (1.93\% of data) shows classic imbalanced learning problems. The wide range of class-specific accuracies (0-87.87\%) suggests the models have not learned generalizable seizure detection patterns but rather dataset-specific biases.


\textbf{Cost-Benefit Analysis:} At 9\,s windows, training is 60\% slower than 5\,s windows, with high variance (9.4\% standard deviation) and marginal absolute performance (14.2\% recall). The substantial computational investment gives minimal improvement on the most critical minority class, indicating future work should prioritize class-imbalance mitigation strategies rather than architectural complexity.

\textbf{Variability Concerns:} The high inter-run variance (coefficient of variation 6-10\%) raises serious concerns about deployment reliability. Clinical monitoring systems require consistent, predictable performance, but the substantial variance suggests model behavior may vary significantly depending on initialization and training dynamics.

\section{Study Limitations}
\label{sec:limitations}

This study has several fundamental limitations that constrain interpretation and applicability:

\subsection{Methodological Limitations}

\textbf{Sample Size:} Five runs per configuration provide insufficient statistical power for robust conclusions, particularly given the observed variance levels. Power analysis suggests 20-30 runs would be needed to detect 2-3 percentage point differences with 80\% power.

\textbf{Single Dataset:} Evaluation on only the \ac{TUEV} dataset limits generalizability. The specific characteristics—particularly severe class imbalance (66.77\% majority class, 1.93\% rarest class)—heavily influence results. Cross-dataset validation is essential before drawing broad conclusions.

\textbf{Hyperparameter Optimization:} Using hyperparameters optimized for the baseline architecture across all variants may underestimate \ac{xLSTM} potential. However, exhaustive hyperparameter search would require much more computational resources than available.

\subsection{Architectural Limitations}

\textbf{Limited Transfer from NLP:} The modest improvements compared to dramatic \ac{xLSTM} success in natural language processing indicate biosignal sequence modeling may require domain-specific architectural innovations rather than adaptations from other fields.

\textbf{Temporal Window Constraints:} The evaluation was restricted to 5, 7, and 9 seconds, which may not capture the full range of temporal contexts relevant for different neurological events. Finer-grained window exploration could better characterize the U-shaped and inverted-U patterns observed.

\subsection{Dataset and Evaluation Limitations}

\textbf{Class Distribution:} The severe class imbalance observed suggests performance would likely vary significantly across datasets with different class distributions.

\textbf{Computational Constraints:} Limited computational resources prevented exploration of larger model configurations, extended training regimes, and comprehensive hyperparameter sweeps.

\section{Summary}
\label{sec:discussion_summary}

This investigation reveals both the potential and significant limitations of integrating \ac{xLSTM} architectures within the \ac{BIOT} framework for \ac{EEG} event classification. While \ac{xLSTM} variants achieve window-specific improvements—\ac{sLSTM} at 5\,s and 9\,s windows, \ac{mLSTM} at 7\,s windows—the lack of statistical significance for balanced accuracy and critical class-specific failures highlight fundamental challenges that extend beyond architectural innovation. The window-dependent performance patterns, where \ac{sLSTM} underperforms at 7\,s and \ac{mLSTM} degrades at 9\,s, show that successful approaches from natural language processing do not directly transfer to biosignal analysis. Only the combined architecture provides consistent stability across windows, though it never achieves best-in-class performance.

Most importantly, this work provides a realistic assessment of current limitations and emphasizes the need for honest reporting of window-specific rather than universal improvements, focus on fundamental challenges like class imbalance, larger-scale evaluation with robust statistical analysis, and recognition that incremental architectural improvements alone may be insufficient for clinical impact. Future research should prioritize addressing these fundamental limitations rather than pursuing incremental architectural modifications.

%% 
%%%%%%%%%%%%%%%%%%%%%%%%%%%%%%%%%%%%%%%%%%%%%%%%%%%%%%%%%%%%%%%%%%%%%%%%%%%%%%%%
